\documentclass[11pt,a4paper]{article}

% --- PAQUETS DE BASE ---
\usepackage[utf8]{inputenc}
\usepackage[T1]{fontenc}
\usepackage[french]{babel}
\usepackage{amsmath, amsfonts, amssymb}
\usepackage{geometry}
\geometry{top=2cm, bottom=2.5cm, left=2cm, right=2cm}

% --- DESIGN & COULEURS ---
\usepackage{xcolor}
\definecolor{BrandPrimary}{HTML}{0D9488}
\definecolor{BrandDark}{HTML}{111827}
\definecolor{BrandLight}{HTML}{F3F4F6}

\usepackage{tcolorbox}
\tcbuselibrary{skins, breakable}

\usepackage{titlesec}
\titleformat{\section}{\color{BrandPrimary}\normalfont\Large\bfseries}{\thesection}{1em}{}[{\titlerule[0.8pt]}]

% --- POLICE ---
\usepackage{helvet}
\renewcommand{\familydefault}{\sfdefault}

\newcommand{\makebrandtitle}{
    \begin{tcolorbox}[enhanced, colback=BrandDark, colframe=BrandPrimary, arc=10pt, boxrule=2pt, halign=center]
        \vspace{0.3cm}
        {\Huge \bfseries \color{white} ZENITH LEARN} \\
        \vspace{0.2cm}
        {\Large \color{BrandPrimary} IA Éthique : Modélisation TWIST 01 (V2 - Strict Math)} \\
        \vspace{0.3cm}
    \end{tcolorbox}
    \vspace{0.5cm}
}

\begin{document}

\makebrandtitle

\section{Analyse Mathématique de l'Engagement}
Le TWIST 01 exige de séparer l'engagement de la complétion. Nous définissons le score d'engagement $S_E$ comme une combinaison pondérée de deux signaux distincts.

\subsection{Formule de l'Engagement}
L'obtention de la formule repose sur l'équilibrage entre le ressenti explicite (la note) et l'effort implicite (le temps passé).
\begin{equation}
S_{E}(u, c) = w_{rat} \cdot \left( \frac{R_{u,c} - R_{min}}{R_{max} - R_{min}} \right) + w_{time} \cdot \min\left( \frac{T_{acc}}{T_{theo}}, \delta \right) \cdot \frac{1}{\delta}
\end{equation}

\subsection{Détail des termes et Rationale}
\begin{itemize}
    \item \textbf{$w_{rat}$ (Poids du Rating) :} Fixé à 0.6. Il représente l'autorité du jugement conscient de l'apprenant.
    \item \textbf{$R_{u,c}$ (Rating brut) :} Note donnée sur une échelle de 1 à 5.
    \item \textbf{$\frac{R_{u,c} - R_{min}}{R_{max} - R_{min}}$ (Normalisation Min-Max) :} Transformation linéaire pour ramener le rating sur $[0,1]$. Indispensable pour éviter que les échelles de notes ne dominent les autres métriques.
    \item \textbf{$w_{time}$ (Poids du temps) :} Fixé à 0.4. Il représente la "preuve de travail" comportementale.
    \item \textbf{$T_{acc}$ (Temps accompli) :} Durée réelle passée par l'apprenant sur les modules du cours $c$.
    \item \textbf{$T_{theo}$ (Temps théorique) :} Durée estimée du cours. Le ratio $\frac{T_{acc}}{T_{theo}}$ mesure l'assiduité relative.
    \item \textbf{$\delta$ (Seuil de saturation) :} Fixé à 1.5. C'est le terme \textbf{anti-spectateur}. Si un utilisateur passe 3 fois plus de temps que prévu, cela n'augmente plus son score d'engagement. On considère que l'utilité marginale décroît violemment au-delà de 150\% du temps prévu.
\end{itemize}

\section{Modélisation de la Complétion}
La complétion $S_C$ est traitée comme un scalaire pur à cette étape.
\begin{equation}
S_{C}(u, c) = \text{Ratio de progression dans le graphe de cours} \in [0, 1]
\end{equation}

\section{Conception du Dataset}
\begin{tabular}{|l|p{10cm}|}
\hline
\rowcolor{BrandDark} \color{white} Colonne & \rowcolor{BrandDark} \color{white} Impact Mathématique \\ \hline
\texttt{rating} & Source de la satisfaction explicite. \\ \hline
\texttt{time\_spent} & Source de l'effort comportemental. \\ \hline
\texttt{completion} & Variable indépendante décorrélée de l'engagement (Principe T01). \\ \hline
\end{tabular}

\end{document}
